\chapter{Application du produit scalaire}
\thispagestyle{fancy}

    \section{Produit scalaire et norme}
    
    \subsection{Propriétés}

    % Soit un vecteur $\vv{u},$ on a :
    % $\vv{u} \cdot \vv{u}=\|\vv{u}\| \times\|\vv{u}\| \times \cos (\vv{u} ; \vv{u})=\|\vv{u}\|^{2} \times \cos 0=\|\vv{u}\|^{2}$\\
    % et\\
    % $\vv{u}\cdot \vv{u}=\vv{u}^{2}$\\
    % On a ainsi $: \vv{u}^{2}=\vv{u} \cdot \vv{u}=\|\vv{u}\|^{2}$
    
    \begin{bprop}
        Propriété : Soit $\vv{u}$ et $\vv{v}$ deux vecteurs. On a :
        \begin{itemDef}
            \item $\vv{u} \cdot \vv{v}=\dfrac{1}{2}\left(\|\vv{u}\|^{2}+\|\vv{v}\|^{2}-\|\vv{u}-\vv{v}\|^{2}\right)$
            \item $\vv{u} \cdot \vv{v}=\dfrac{1}{2}\left(\|\vv{u}+\vv{v}\|^{2}-\|\vv{u}\|^{2}-\|\vv{v}\|^{2}\right)$
            \item $\vv{u} \cdot \vv{v}=\dfrac{1}{4}\left(\|\vv{u}+\vv{v}\|^{2}-\|\vv{u}-\vv{v}\|^{2}\right)$
        \end{itemDef}
    \end{bprop}
    
    \begin{bdemot}{Démonstration de la première formule}
        \begin{align*}
            \|\vv{u}-\vv{v}\|^{2}&=(\vv{u}-\vv{v})^{2} \\
            &=\vv{u}^{2}-2 \vv{u} \cdot \vv{v}+\vv{v}^{2} \\
            &=\|\vv{u}\|^{2}-2 \vv{u} \cdot \vv{v}+\|\vv{v}\|^{2}
        \end{align*}
        Donc $\vv{u} \cdot \vv{v}=\dfrac{1}{2}\left(\|\vv{u}\|^{2}+\|\vv{u}\|^{2}-\|\vv{u}-\vv{v}\|^{2}\right)$
    \end{bdemot}

    \begin{bprop}
        Soit $A, B$ et $C$ trois points du plan. On a :
        $\vv{A B} \cdot \vv{A C}=\dfrac{1}{2}\left(A B^{2}+A C^{2}-B C^{2}\right)$
    \end{bprop}
    
    \begin{bdemo}
        \begin{align*}
            \vv{A B} \cdot \vv{A C}&=\dfrac{1}{2}\left(\|\vv{A B}\|^{2}+\|\vv{A C}\|^{2}-\|\vv{A B}-\vv{A C}\|^{2}\right) \\
            &=\dfrac{1}{2}\left(A B^{2}+A C^{2}-\|\vv{C B}\|^{2}\right) \\
            &=\dfrac{1}{2}\left(A B^{2}+A C^{2}-B C^{2}\right)
        \end{align*}
    \end{bdemo}
    
    \begin{bexemple}
        Si $\mathrm{AB}=6, \mathrm{AC}=5$ et $\mathrm{CB}=4$ alors:
        $$
        \vv{A B} \cdot \vv{A C}=\frac{1}{2}\left(A B^{2}+A C^{2}-C B^{2}\right)=\frac{1}{2}\left(6^{2}+5^{2}-4^{2}\right)=\frac{1}{2}(36+25-16)=\frac{45}{2}
        $$
    \end{bexemple}
    
    \begin{bmethodet}{Calculer un produit scalaire à l'aide des normes}
        \grasInMet{Enoncé :}\\
        On considère la figure ci-dessous, calculer le produit scalaire $\vv{C G} \cdot \vv{C F}$
        \begin{center}
            \includegraphics[scale=0.8]{1.04.IMG-02}
        \end{center}
        \ \\
        \grasInMet{Réponse :}\\
        $\vv{C G} \cdot \vv{C F}=\frac{1}{2}\left(C G^{2}+C F^{2}-G F^{2}\right)$\\
        $=\dfrac{1}{2}\left(6^{2}+7^{2}-3^{2}\right)$\\
        $=38$
    \end{bmethodet}
    
    \subsection{Théorème d'Al Kashi}
    
    \begin{bthm}
        Dans un triangle $\mathrm{ABC},$ on a, avec les notations de la figure :
        $$
        a^{2}=b^{2}+c^{2}-2 b c \cos \hat{A}
        $$
        \begin{center}
            \includegraphics[scale=0.8]{1.04.IMG-03}
        \end{center}
    \end{bthm}
    
    \begin{bdemo}
        $\vv{A B} \cdot \vv{A C}=A B \times A C \times \cos \hat{A}=b c \cos \hat{A}$\\
        et\\
        $\vv{A B} \cdot \vv{A C}=\frac{1}{2}\left(A B^{2}+A C^{2}-B C^{2}\right)=\frac{1}{2}\left(b^{2}+c^{2}-a^{2}\right)$\\
        Donc $: \frac{1}{2}\left(b^{2}+c^{2}-a^{2}\right)=b c \cos \hat{A}$\\
        Soit $: b^{2}+c^{2}-a^{2}=2 b c \cos \hat{A}$\\
        Soit encore $: a^{2}=b^{2}+c^{2}-2 b c \cos \hat{A}$\\
    \end{bdemo}
    
    \begin{bmethodet}{Calculer un angle à l'aide des normes}
        \grasInMet{Enoncé :}\\
        On considère la figure ci-dessous, calculer la mesure de l'angle $\widehat{B A C}$ au degré près.
        \begin{center}
            \includegraphics[scale=0.8]{1.04.IMG-05}
        \end{center}
        \ \\
        \grasInMet{Réponse :}\\
        D'après le théorème d'Al Kashi, on a :\\
        $C B^{2}=A B^{2}+A C^{2}-2 \times A B \times A C \times \cos \widehat{B A C}$\\
        $4^{2}=6^{2}+5^{2}-2 \times 6 \times 5 \times \cos \widehat{B A C}$\\
        $16=36+25-60 \cos \widehat{B A C}$\\
        $60 \cos \widehat{B A C}=36+25-16$\\
        $60 \cos \widehat{B A C}=45$\\
        $\cos \widehat{B A C}=\dfrac{45}{60}$\\
        $\cos \widehat{B A C}=\dfrac{3}{4}$\\
        $\widehat{B A C} \approx 41^{\circ}$
    \end{bmethodet}
    
    \section{Vecteur Normal}
    
    \subsection{définition et propriétés}
    
    \begin{bdefi}
        Un \grasInDef{vecteur normal} à une droite $d$ quelconque du plan est un vecteur non nul et orthogonal à un vecteur directeur de $d$.
    \end{bdefi}
    
    \begin{bprop}
        Deux droites du plan sont perpendiculaires si et seulement si un vecteur normal de l'une est orthogonal à un vecteur normal de l'autre.
    \end{bprop}
    
    \begin{bdemo}
        On suppose que $d$ et $d^{\prime}$ sont perpendiculaires. Si $\vv{u}$ est un vecteur directeur de $d$ et $\vv{v}$ de $d^{\prime},$ alors $\vv{u}$ et $\vv{v}$ sont orthogonaux. $\vv{v}$ étant normal à $d$ et $\vv{u}$ à $d^{\prime},$ la propriété est vérifiée.\\
        Réciproquement, si $\vv{u},$ vecteur normal à $d,$ est orthogonal à $\vv{v},$ vecteur normal à $d^{\prime},$ alors $\vv{v}$ est un vecteur directeur de $d$ et $\vv{u}$ de $d^{\prime} .$ Ayant des vecteurs directeurs orthogonaux, d et $d^{\prime}$ sont perpendiculaires.
    \end{bdemo}
    
    \begin{bmethode}
        \grasInMet{Enoncé:}\\
        Soit $d$ une droite de vecteur directeur $\Coord{u}{3}{5}$.
        \begin{enumMet}
            \item Déterminer une condition sur les coordonnées d'un vecteur $\vv{n}$ non nul pour qu'il soit normal à $d$.
            \item Le vecteur $\Coord{v}{-10}{6}$ est-il un vecteur normal à $d$?
        \end{enumMet}
        
        \grasInMet{Réponse:}\\
        \begin{enumMet}
            \item Si $\Coord{n}{x}{y}$ est un vecteur normal à $d$, alors il est orthogonal à $\vv{u}$.\\
            Donc $\vv{u} \cdot \vv{n}=3 x+5 y=0$ soit $y=-\dfrac{3}{5} x$.\\
            Les coordonnées de $\vv{n}$ sont donc de la forme $\CoordSV{x}{-\dfrac{3}{5}x}$ avec $x \in \R$.
            \item On calcule le produit scalaire $\vv{u}\cdot\vv{v}$\\
            $\vv{u}\cdot\vv{v}=3\times(-10)+5\times6=-30+30=0$\\
            le vecteur $\vv{v}$ est orthogonal au vecteur directeur $\vv{u}$ de la droite $d$, donc $\vv{v}$ est un vecteur normal à $d$. 
        \end{enumMet}
    \end{bmethode}
    
    \subsection{Équations cartésiennes et vecteur normal}
    
    \begin{brappel}
        Toute droite $D$ admet une équation de la forme $a x+b y+c=0$ avec $(a ; b) \neq(0 ; 0) .$ Un \grasInDef{vecteur directeur} de $D$ est $\Coord{u}{-b}{a}$ \\
        Cette équation est appelée \grasInDef{équation cartésienne} de la droite $D$.
    \end{brappel}
        
    \begin{bprop}
        Soient $a, b$ et $c$ trois réels tels que $a$ et $b$ ne sont pas simultanément nuls.\\
        La droite d'équation cartésienne $a x+b y+c=0$ admet pour \grasInDef{vecteur normal} le vecteur $\Coord{n}{a}{b}$.\\
        \ \\
        Réciproquement, toute droite ayant pour \grasInDef{vecteur normal} le vecteur non nul $\Coord{n}{a}{b}$ admet une équation cartésienne de la forme $a x+b y+c=0$.
    \end{bprop}
        
    \begin{bexemple}
        La droite d'équation cartésienne $3 x-4 y+5=0$ admet pour vecteur normal le vecteur $\Coord{n}{3}{-4}$  et pour vecteur directeur le vecteur $\Coord{u}{4}{3}$.\\
        On a bien $\vv{n} \cdot \vv{u}=0$
    \end{bexemple}
    
    \begin{bmethode}
        \grasInMet{Enoncé:}\\
        Dans un repère orthonormé, déterminer, de deux façons différentes, une équation cartésienne de la droite $d$ passant par le point $\mathrm{A}(5 ;-1)$ et de vecteur normal $\Coord{n}{2}{-3}$.
        
        \grasInMet{Réponse:}\\
        \begin{enumMet}
            \item 1. On lit sur le vecteur normal que $a=2$ et $b=-3$.
            Donc une équation de $d$ est de la forme $2 x-3 y+c=0$. $\mathrm{A}(5 ;-1) \in d$ donc $2 \times 5-3 \times(-1)+c=0$ et donc $c=-13$.\\
            Ainsi, une équation de la droite $d$ est $2 x-3 y-13=0$.
            \item Soit $\mathrm{M}(x ;y)$ appartenant à $d .$ Alors $\vv{\mathrm{AM}}$ est un vecteur directeur de $d$ et est orthogonal au vecteur $\vv{n}$.\\
            $\vv{\mathrm{AM}} \cdot \vv{n}=0 \iff \CoordSV{x-5}{y+1} \cdot \CoordSV{2}{-3}=0 \iff 2(x-5)-3(y+1)=0 \iff 2 x-3 y-13=0$\\
            Donc une équation de la droite $d$ est $2 x-3 y-13=0$
        \end{enumMet}
    \end{bmethode}

